\section{Objetivos}
\subsection{Objetivo General}

\begin{itemize}
    \item Estudiar el desempeño operativo de sistemas formados por transformadores monofásicos acoplados en paralelo, examinando de manera integral los parámetros técnicos que rigen su comportamiento tanto en condición de vacío como ante la presencia de cargas.
\end{itemize}

\subsection{Objetivos Específicos}

\begin{itemize}
    \item Comprobar los requisitos fundamentales de polaridad, relación de transformación y equivalencia de impedancias previos al acoplamiento en paralelo de las unidades monofásicas.
    
    \item Medir y analizar el fenómeno de la corriente circulatoria generado por variaciones en las relaciones de transformación de los equipos.
    
    \item Analizar el reparto de corriente y potencia entre las distintas unidades acopladas, considerando la influencia directa de las impedancias de cortocircuito del circuito equivalente.
    
    \item Contrastar los resultados experimentales obtenidos en el laboratorio con los valores teóricos previstos para la corriente circulatoria, tomando en cuenta las posibles desviaciones técnicas.
    
\end{itemize}